\begin{lemma}[Nakayama's lemma]
\label{lemma-NAK}
\begin{reference}
\cite[1.M Lemma (NAK) page 11]{MatCA}
\end{reference}
\begin{history}
We quote from \cite{MatCA}: ``This simple but
important lemma is due to T.~Nakayama, G.~Azumaya and W.~Krull. Priority
is obscure, and although it is usually called the Lemma of Nakayama, late
Prof.~Nakayama did not like the name.''
\end{history}
Let $R$ be a ring with Jacobson radical $\text{rad}(R)$.
Let $M$ be an $R$-module. Let $I \subset R$
be an ideal.
\begin{enumerate}
\item
\label{item-nakayama}
If $IM = M$ and $M$ is finite, then there exists an $f \in 1 + I$ such that
$fM = 0$.
\item If $IM = M$, $M$ is finite, and $I \subset \text{rad}(R)$, then $M = 0$.
\item If $N, N' \subset M$, $M = N + IN'$, and $N'$ is finite,
then there exists an $f \in 1 + I$ such that $fM \subset N$ and $M_f = N_f$.
\item If $N, N' \subset M$, $M = N + IN'$, $N'$ is finite, and
$I \subset \text{rad}(R)$, then $M = N$.
\item If $N \to M$ is a module map, $N/IN \to M/IM$ is
surjective, and $M$ is finite, then there exists an $f \in 1 + I$
such that $N_f \to M_f$ is surjective.
\item If $N \to M$ is a module map, $N/IN \to M/IM$ is
surjective, $M$ is finite, and $I \subset \text{rad}(R)$,
then $N \to M$ is surjective.
\item If $x_1, \ldots, x_n \in M$ generate $M/IM$ and $M$ is finite,
then there exists an $f \in 1 + I$ such that $x_1, \ldots, x_n$
generate $M_f$ over $R_f$.
\item If $x_1, \ldots, x_n \in M$ generate $M/IM$, $M$ is finite, and
$I \subset \text{rad}(R)$, then $M$ is generated by $x_1, \ldots, x_n$.
\item If $IM = M$, $I$ is nilpotent, then $M = 0$.
\item If $N, N' \subset M$, $M = N + IN'$, and $I$ is nilpotent then $M = N$.
\item If $N \to M$ is a module map, $I$ is nilpotent, and $N/IN \to M/IM$
is surjective, then $N \to M$ is surjective.
\item If $\{x_\alpha\}_{\alpha \in A}$ is a set of elements of $M$
which generate $M/IM$ and $I$ is nilpotent, then $M$ is generated
by the $x_\alpha$.
\end{enumerate}
\end{lemma}

\begin{proof}
Proof of (\ref{item-nakayama}). Choose generators $y_1, \ldots, y_m$ of $M$
over $R$. For each $i$ we can write $y_i = \sum z_{ij} y_j$ with
$z_{ij} \in I$ (since $M = IM$).
In other words $\sum_j (\delta_{ij} - z_{ij})y_j = 0$.
Let $f$ be the determinant of the $m \times m$ matrix
$A = (\delta_{ij} - z_{ij})$. Note that $f \in 1 + I$
(since the matrix $A$ is entrywise congruent to the
$m \times m$ identity matrix modulo $I$).
By Lemma \ref{lemma-matrix-left-inverse} (1),
there exists an $m \times m$
matrix $B$ such that $BA = f 1_{m \times m}$. Writing out we see that
$\sum_{i} b_{hi} a_{ij} = f \delta_{hj}$ for all
$h$ and $j$; hence, $\sum_{i, j} b_{hi} a_{ij} y_j
= \sum_{j} f \delta_{hj} y_j = f y_h$ for every $h$.
In other words, $0 = f y_h$ for every $h$ (since each
$i$ satisfies $\sum_j a_{ij} y_j = 0$).
This implies that $f$ annihilates $M$.

\medskip\noindent
By Lemma \ref{lemma-contained-in-radical} an element of $1 + \text{rad}(R)$ is
invertible element of $R$. Hence we see that (\ref{item-nakayama}) implies
(2). We obtain (3) by applying (1) to $M/N$ which is finite as $N'$ is finite.
We obtain (4) by applying (2) to $M/N$ which is finite as $N'$ is finite.
We obtain (5) by applying (3) to $M$ and the submodules $\Im(N \to M)$
and $M$. We obtain (6) by applying (4) to $M$ and the submodules
$\Im(N \to M)$ and $M$.
We obtain (7) by applying (5) to the map $R^{\oplus n} \to M$,
$(a_1, \ldots, a_n) \mapsto a_1x_1 + \ldots + a_nx_n$.
We obtain (8) by applying (6) to the map $R^{\oplus n} \to M$,
$(a_1, \ldots, a_n) \mapsto a_1x_1 + \ldots + a_nx_n$.

\medskip\noindent
Part (9) holds because if $M = IM$ then $M = I^nM$ for all $n \geq 0$
and $I$ being nilpotent means $I^n = 0$ for some $n \gg 0$. Parts
(10), (11), and (12) follow from (9) by the arguments used above.
\end{proof}
